\chapter{Metodología}
El desarrollo del proyecto se realizó mediante un proceso por etapas 
orientado a la construcción de un prototipo funcional de videoclip 
gamificado en realidad virtual (VR, por sus siglas en inglés). La 
primera etapa del fue la de la preparación del entorno de desarrollo y la 
implementación de las mecánicas básicas de los minijuegos. Posteriormente 
se incorporaron los mecanismos de interacción en VR, el registro de 
métricas de desempeño, el modelado de estados del jugador y el sistema 
de ajuste dinámico de dificultad (DDA, por sus siglas en inglés). 
Finalmente, los distintos módulos se integraron con el flujo general de 
la experiencia y se realizaron pruebas orientadas a comprobar su 
funcionamiento conjunto.

\section{Enfoque general del desarrollo}
El proyecto se abordó como el desarrollo de un prototipo interactivo en 
tiempo real. Para ello, el sistema se dividió en componentes responsables 
de tareas específicas. Las tareas fueron: interacción en VR, lógica de 
los minijuegos, captura de métricas, representación del estado del 
jugador, adaptación dinámica y comunicación con el flujo global de la 
experiencia. Esta separación permitió desarrollar y validar cada 
componente antes de integrarlo con los demás. 

El sistema de interacción generó los eventos correspondientes a las 
acciones realizadas por el usuario. Dichos eventos alimentaron el 
sistema de métricas, y las métricas registradas se utilizaron para 
determinar el comportamiento del sistema adaptativo. De esta manera, se 
estableció un flujo general de interacción medición, interpretación y 
adaptación. 

La adaptación se diseño utilizando un enfoque basado en reglas. Esta 
decisión permitió relacionar de manera explícita las variables 
observadas durante el gameplay con los cambios realizados sobre los 
minijuegos, facilitando la inspección y ajuste de las condiciones 
utilizadas por el sistema. Así se construyo el modelo sin que fuera un 
modelo clínico o psicológico del usuario, sino una representación 
operativa de su comportamiento basado en información generada 
directamente durante la ejecución.

\section{Preparación del entono de desarrollo en VR}
Lo primero fue configurar el proyecto en Unity para soportar interacción 
XR y se habilitaron los componentes necesarios para reconocer la 
posición y orientación del visor Meta Quest y los dispositivos de 
entrada. Luego, se configuró la representación del usuario dentro de las 
escenas de manera que los movimientos realizados con el dispositivo 
pudieran trasladarse al entorno virtual. A partir de esta configuración 
se comprobó el funcionamiento de la cámara, los controles y las 
interacciones básicas antes de incorporar las mecánicas específicas de 
los minijuegos. 

También se estableció una estructura común para los elementos que 
posteriormente serían utilizados por los distintos sistemas. El objetivo 
fue evitar que cada minijuego tuviera que implementar de forma 
independiente mecanismos equivalentes de entrada, detección o 
comunicación.

Los recursos tridimensionales necesarios para los minijuegos se 
obtuvieron, principalmente, de la \textit{Asset Store} de Unity, pero 
recursos específicos como el modelo de los músicos fueron creados 
utilizando Blender. Todos los modelos siguen una misma dirección visual 
estilizada y un estilo \textit{low-poly}, procurando mantener una 
complejidad adecuada para la ejecución de la experiencia en el Meta 
Quest 2. Una vez preparados o recolectados, los recursos fueron 
importados a Unity y configurados con los componentes necesarios para 
permitir su utilización dentro de las escenas.

\section{Diseño e implementación de los minijuegos}
Se desarrollaron tres minijuegos para representar diferentes formas de 
interacción dentro del videoclip gamificado. Cada uno fue diseñado 
alrededor de una actividad breve que pudiera realizarse utilizando 
acciones físicas dentro del entorno virtual y que permitiera generar 
eventos cuantificables para el sistema adaptativo.

Antes de implementar cada minijuego se definieron sus objetivos, las 
acciones válidas del jugador, las condiciones de acierto y error y los 
elementos susceptibles de modificación durante la ejecución. Esta 
definición permitió separar las reglas fundamentales de cada actividad 
de los parámetros que posteriormente podrían alterarse mediante el DDA.

Las mecánicas se implementaron inicialmente sin adaptación dinámica. 
Esta primera versión permitió comprobar que las acciones pudieran 
completarse correctamente y que existiera una respuesta clara ante los 
resultados del jugador. Posteriormente, se incorporaron los puntos de 
medición necesarios para registrar el desempeño y los parámetros que 
podían ser modificados por el sistema adaptativo.

Aunque cada minijuego utilizó mecánicas diferentes, se mantuvo una 
estructura común para comunicar eventos de interacción y desempeño. Esto 
permitió que el sistema de métricas y adaptación operara sobre 
información equivalente sin depender directamente de la lógica interna de 
una mecánica específica.

\section{Implementación del sistema base de interacción}
Una vez establecidas las mecánicas principales, se implementó el sistema 
encargado de reconocer y procesar las acciones realizadas por el jugador 
dentro del entorno virtual. Este sistema recibió la información 
proveniente de los dispositivos de entrada y la relacionó con los 
objetos interactivos presentes en cada minijuego. La detección de 
interacción se realizó utilizando los mecanismos de colisión y detección 
disponibles en Unity. 

Los objetos relevantes fueron configurados para identificar eventos como 
contacto, selección, impacto o activación, dependiendo de las necesidades 
de cada mecánica. A partir de estos eventos se ejecutaron las reglas 
correspondientes para determinar si una acción representaba un acierto, 
un error o una interacción que no debía ser contabilizada. También se 
implementaron condiciones específicas para controlar cuándo una acción 
podía ser considerada válida. Esto permitió reducir registros duplicados 
y evitar que contactos involuntarios fueran interpretados como múltiples 
interacciones.

Cada acción relevante generó retroalimentación inmediata para comunicar 
al jugador el resultado de su interacción. Se utilizaron respuestas 
visuales, sonoras y hápticas según las características de cada minijuego. 
Estos mecanismos se integraron con la propia lógica de interacción para 
que la respuesta ocurriera en el momento en que el sistema reconocía el 
evento.

\section{Diseño y registro de métricas de desempeño}
Después de comprobar el funcionamiento de las interacciones, se 
implementó un sistema de telemetría destinado a registrar información del 
desempeño del jugador durante la ejecución. Para ello, se seleccionaron 
como variables principales los aciertos, errores, tiempo de reacción, 
precisión y frecuencia de interacción. Estos datos permitieron analizar 
de manera más objetiva la respuesta del jugador ante las distintas 
actividades y obtener información útil para evaluar el funcionamiento de 
la experiencia. 

Los aciertos y errores se registraron a partir del resultado producido 
por las reglas de cada minijuego. Para evitar los que los valores 
dependieran de la lógica interna de una actividad específica, cada 
minijuego comunicó al sistema de métricas el resultado de las 
interacciones relevantes mediante eventos comunes. El tiempo de reacción 
se calculó utilizando el intervalo entre el momento en que una 
oportunidad de interacción se encontraba disponible para el jugador y el 
instante en que se detectaba una respuesta válida. Esta información 
permitió observar cambios en la rapidez de respuesta durante diferentes 
momentos de la experiencia. 

La precisión se utilizó para representar qué tan correctamente se 
ejecutaban aquellas acciones cuya mecánica permitía evaluar la 
correspondencia entre el movimiento realizado y el objetivo establecido. 
Debido a que los minijuegos utilizaban formas de interacción diferentes, 
esta métrica fue normalizada de acuerdo con las características de cada 
actividad antes de ser utilizada por el sistema de adaptación. La 
frecuencia de interacción se obtuvo a partir de la cantidad de acciones 
relevantes detectadas durante un intervalo de tiempo. Esta variable 
permitió diferenciar períodos de actividad continua de momentos en los 
que el usuario reducía o detenía su participación.

Las métricas se mantuvieron disponibles durante la ejecución de la 
experiencia para que pudieran ser consultadas por los sistemas de estados 
y adaptación. Además de los valores acumulados, se consideraron 
observaciones recientes del comportamiento para evitar que el desempeño 
de los primeros momentos de la experiencia condicionara permanentemente 
las decisiones posteriores. De esta manera, el sistema podía actualizar 
la interpretación del desempeño del jugador conforme avanzaba la sesión 
y respondía a las distintas actividades. Esto permitió que las decisiones 
de adaptación se basaran tanto en la evolución general de los resultados 
como en el comportamiento observado en los momentos más recientes.

\section{Modelado de estados del jugador}
A partir de las métricas registradas se desarrolló una representación 
simplificada del estado de interacción del jugador. Se definieron cuatro 
estados operativos: quieto, interactuando, sobrecargado y transicionando. 
Estos estados no representaron diagnósticos psicológicos, sino categorías 
funcionales para la experiencia que resumían patrones observables durante 
su ejecución.

El estado “quieto” se utilizó cuando el sistema detectaba una reducción 
sostenida de las interacciones esperadas. El estado “interactuando” 
representó una participación regular dentro de los parámetros esperados 
para la actividad. El estado “sobrecargado” se usó cuando la combinación 
de métricas mostraba dificultades persistentes para responder 
adecuadamente a las exigencias actuales del minijuego. Finalmente, el 
estado “transicionando” se utilizó durante períodos en los que el flujo 
de la experiencia cambiaba de una actividad a otra y las métricas no 
debían ser interpretadas utilizando las mismas condiciones de una 
interacción activa.

La clasificación no se realizó a partir de un único evento. Se analizaron 
varias observaciones recientes para reducir la posibilidad de cambiar de 
estado como consecuencia de un error aislado o una pausa breve. Asimismo, 
se establecieron condiciones para controlar las transiciones entre 
estados y evitar cambios continuos cuando las métricas se encontraban 
cerca de un límite. 

El estado obtenido se mantuvo disponible como una entrada adicional para 
el sistema adaptativo, complementando las métricas individuales. De esta 
manera, las decisiones de ajuste podían considerar tanto valores 
concretos de desempeño como tendencia general identificada durante el 
período reciente de interacción. Esto permitía que el sistema no 
reaccionara únicamente a resultados aislados, sino que también tuviera en 
cuenta la evolución del comportamiento del jugador a lo largo del tiempo.

\section{Implementación del DDA}
El sistema DDA se desarrolló después de disponer de interacciones 
estables, métricas registrables y una representación del estado del 
jugador. Su funcionamiento se basó en reglas que relacionaban tendencias 
de desempeño con modificaciones sobre parámetros controlables de cada 
minijuego. Estas reglas permitían adaptar progresivamente la experiencia 
según el rendimiento observado, evitando que la dificultad permaneciera 
fija durante toda la sesión.

Se definieron tres categorías principales de parámetros adaptativos: 
velocidad, cantidad de objetos y complejidad de la tarea. La aplicación 
concreta de cada parámetro dependió de las características del minijuego. 
De esta manera, el sistema adaptativo determinaba la dirección del 
cambio, mientras que cada mecánica establecía como debía materializarse 
dicho ajuste.

Las decisiones no se realizaron después de cada interacción individual. 
El sistema acumuló información durante períodos de observación y evaluó 
tendencias antes de aplicar una modificación. Cuando el desempeño se 
mantenía por debajo de los valores esperados o el jugador permanecía en 
el estado de sobrecarga, se aplicaban ajustes orientados a reducir las 
exigencias de la actividad. Cuando se observaba un desempeño estable y 
sostenido, el sistema podía aumentar gradualmente la exigencia. 

También se definieron situaciones en las que la dificultad debía 
mantenerse. Esto evitó que pequeñas variaciones produjeran modificaciones 
innecesarias. Después de ejecutar un cambio se permitió un período de 
estabilización antes de realizar una nueva evaluación, con el propósito 
de observar el efecto de la modificación y reducir oscilaciones continuas 
entre niveles dificultad. 

Los ajustes se limitaron a rangos previamente definidos para cada 
minijuego. De esta manera, el sistema no podía reducir o aumentar los 
parámetros hasta valores que alteren las reglas fundamentales de la 
actividad o produjeran condiciones incompatibles con la experiencia 
musical. El DDA se mantuvo separado de la implementación específica de 
los minijuegos mediante eventos y parámetros configurables. Esta 
separación permitió utilizar una lógica general de adaptación mientras 
cada actividad conservaba el control sobre la forma en que una 
modificación afectaba su \textit{gameplay}.

\section{Integración del sistema adaptativo con la experiencia global}
Una vez implementados los componentes anteriores, el sistema se integró 
con el flujo general del videoclip gamificado. La integración se realizó 
mediante eventos que permitieron recibir información procedente del 
sistema global y comunicar resultados relacionados con el desempeño del 
jugador. De esta forma, las acciones realizadas durante los minijuegos 
podían influir en el desarrollo de la experiencia y en la actualización 
de las métricas correspondientes.

El sistema adaptativo recibió notificaciones necesarias para conocer el 
contexto general de ejecución, por ejemplo, el inicio o finalización de 
una actividad o el ingreso a un período de transición. Estas señales se 
utilizaron para controlar cuándo debía comenzar o detenerse la captura e 
interpretación de las métricas. De manera inversa, los minijuegos y los 
sistemas de interacción emitieron información relacionada con aciertos, 
errores, progreso y desempeño. Esta comunicación permitió mantener 
desacoplada la lógica adaptativa de los componentes responsables del 
flujo global de la experiencia.

\section{Pruebas de calibración y ajustes}
Las pruebas del sistema adaptativo se realizaron provocando diferentes 
patrones de desempeño para comprobar que las reglas reaccionaran en la 
dirección esperada. Se evaluaron escenarios de desempeño sostenido, 
sucesión de errores, ausencia temporal de interacción y recuperación 
posterior del jugador. Esto permitió observar las transiciones entre 
estados y verificar que los parámetros de dificultad fueran modificados 
únicamente cuando se cumplían las condiciones establecidas.

Posteriormente, se ajustaron los umbrales, períodos de observación, 
límites de los parámetros y tiempos entre modificaciones. Esta 
calibración se hizo para evitar oscilaciones frecuentes entre 
configuraciones de dificultad. Con esto el sistema ya solo aplicaba una 
modificación cuando existían suficientes indicios de que el nivel de 
dificultad debía mantenerse, aumentarse o reducirse. Finalmente, el 
prototipo fue probado como una experiencia integrada utilizando el Meta 
Quest 2 para comprobar la interacción física, la retroalimentación, la 
captura de datos y la respuesta del sistema adaptativo. 




